% =====================================================================
\chapter*{Summary}
\addcontentsline{toc}{chapter}{Summary and References}
\markboth{Summary}{}
% =====================================================================

Table~\ref{tab:summary} collects the main result of every lab. The numbers are not comparable across
rows, since each lab uses a different dataset and task; they show what each algorithm achieved on its
own problem.

\begin{table}[H]
\centering\small
\caption{Main results of the thirteen labs (all on held-out test data).}
\label{tab:summary}
\begin{tabularx}{\linewidth}{@{}r l >{\raggedright\arraybackslash}X >{\raggedright\arraybackslash}X@{}}
\toprule
\textbf{Lab} & \textbf{Algorithm} & \textbf{Dataset} & \textbf{Main test result}\\
\midrule
1  & k-NN                   & Wine                 & Accuracy 92.5\%, macro F1 0.924\\
2  & Linear / polynomial regression & California Housing & $R^2=0.621$, RMSE 0.712 lakh USD (degree 3)\\
3  & Gradient descent       & Concrete strength    & $R^2=0.835$, RMSE 6.66 MPa (same as exact solution)\\
4  & MLE vs MAP             & Diabetes progression & RMSE 59.6 (MAP) vs 63.3 (MLE)\\
5  & Naive Bayes            & SMS spam             & Accuracy 99.0\%, spam F1 0.963, AUC 0.991\\
6  & Logistic regression    & Breast cancer        & Accuracy 97.3\%, recall 97.6\%, AUC 0.997\\
7  & Perceptron             & Iris (two species)   & Accuracy 100\%, F1 1.000\\
8  & Multi-layer perceptron & Iris                 & Accuracy 96.7\%, macro F1 0.966\\
9  & RBF network            & Wheat seeds          & Accuracy 96.8\%, macro F1 0.968\\
10 & SVM (RBF kernel)       & Sonar                & Accuracy 95.1\% (CV 85.0\%), AUC 0.969\\
11 & SVR (RBF kernel)       & Auto MPG             & $R^2=0.887$, RMSE 2.67 mpg\\
12 & Decision tree (ID3)    & Mushroom             & Accuracy 100\%, macro F1 1.000\\
13 & Genetic algorithm      & Breast cancer        & 12 of 30 features; CV accuracy 96.8\% $\to$ 98.5\%\\
\bottomrule
\end{tabularx}
\end{table}

\section*{Lessons that came up again and again}
\begin{itemize}
\item \textbf{Preparing the data mattered as much as the algorithm.} Scaling the features was worth 25
      points for k-NN (Lab~1); taking the log of concrete age (Lab~3) made a simple linear model work
      well; standardising the inputs kept the polynomial models in Lab~2 numerically stable for longer;
      and removing features improved k-NN's cross-validation score on the breast-cancer data (Lab~13).
\item \textbf{Every model needed some way to control its complexity.} This was $k$ in k-NN, the polynomial degree in Lab~2, $\lambda$ in
      MAP, the hidden-layer size of the MLP, the number of centres and the width in the RBF network,
      $C$ and $\gamma$ in the SVM and SVR, and the depth of the tree. In every case I chose it with
      cross-validation on the training data, never with the test set.
\item \textbf{The same ideas link many of the algorithms.} Least squares appears in linear regression, in
      MLE under Gaussian noise and in the output layer of the RBF network. Gradient descent trains logistic
      regression and the MLP. Lagrange multipliers and kernels connect the SVM and SVR. The perceptron's
      failure on XOR is the reason the MLP exists.
\item \textbf{Small test sets give noisy numbers.} In Labs 1 and 10 the test accuracy and the CV
      accuracy differed by 7--10 points. Reporting both, and saying which one to trust, is more honest
      than reporting one number.
\end{itemize}

Writing every algorithm from scratch took more time than calling a library, but it forced me to
understand details that are easy to skip in a lecture: how ties are broken in k-NN, why log-probabilities
are needed in Naive Bayes, what the SMO update actually does, and what happens to an ID3 branch with no
training examples.

\section*{References}
\begin{enumerate}[label={[\arabic*]}, leftmargin=2.2em, itemsep=1pt, font=\small]
\small
\item M.~T. Hagan, H.~B. Demuth, M.~H. Beale and O. De~Jesús, \emph{Neural Network Design}, 2nd ed., 2014 (Chapters 2, 3, 4 and 8).
\item S. Haykin, \emph{Neural Networks and Learning Machines}, 3rd ed., Pearson, 2009.
\item C.~M. Bishop, \emph{Pattern Recognition and Machine Learning}, Springer, 2006.
\item J.~R. Quinlan, ``Induction of decision trees'', \emph{Machine Learning}, 1(1):81--106, 1986.
\item J.~C. Platt, ``Sequential minimal optimization: A fast algorithm for training support vector machines'', Microsoft Research Technical Report MSR-TR-98-14, 1998.
\item R.-E. Fan, P.-H. Chen and C.-J. Lin, ``Working set selection using second order information for training support vector machines'', \emph{Journal of Machine Learning Research}, 6:1889--1918, 2005.
\item A.~J. Smola and B. Schölkopf, ``A tutorial on support vector regression'', \emph{Statistics and Computing}, 14:199--222, 2004.
\item D. Dua and C. Graff, \emph{UCI Machine Learning Repository}, University of California, Irvine. Source of the Wine, Concrete, SMS Spam, Breast Cancer Wisconsin, Iris, Seeds, Sonar, Auto MPG and Mushroom datasets.
\item B. Efron, T. Hastie, I. Johnstone and R. Tibshirani, ``Least angle regression'', \emph{Annals of Statistics}, 32(2):407--499, 2004. Source of the diabetes dataset.
\item I.-C. Yeh, ``Modeling of strength of high-performance concrete using artificial neural networks'', \emph{Cement and Concrete Research}, 28(12):1797--1808, 1998.
\item T.~A. Almeida, J.~M. Gómez Hidalgo and A. Yamakami, ``Contributions to the study of SMS spam filtering'', \emph{Proc. ACM DocEng}, 2011.
\item R.~P. Gorman and T.~J. Sejnowski, ``Analysis of hidden units in a layered network trained to classify sonar targets'', \emph{Neural Networks}, 1:75--89, 1988.
\item M. Bohanec and V. Rajkovič, ``Knowledge acquisition and explanation for multi-attribute decision making'', \emph{8th Intl. Workshop on Expert Systems and their Applications}, Avignon, 1988.
\end{enumerate}
